\documentclass[12pt]{article}
\usepackage[utf8]{inputenc}
\usepackage[T1]{fontenc}
\usepackage{hyperref}
\title{Ulithian Language Resources for K-12 Curriculum Development}
\author{Compiled by Claude Sonnet 5 (claude-sonnet-5) for Andr\'as Kornai}
\date{}
\begin{document}
\maketitle
\section{Introduction}
Ulithian (ISO 639-3 uli) is an Austronesian (Micronesian) language of Ulithi Atoll in Yap State, Federated States of Micronesia. Sources report some 3,000 speakers, of whom about 700 live on the atoll. Wikipedia notes a UNESCO ``severely endangered'' classification.
\section{Primary Educational Resources}
No school textbooks or primers were found.
\section{Linguistic and Reference Materials}
The Habele Outer Island Education Fund published a Ulithian--English and English--Ulithian dictionary in 2010, described as the first comprehensive one: \url{https://www.comfsm.fm/news/releases/habele4.html}. Overview: \url{https://wikipedia.com/wiki/Ulithian}. A grammar by Sohn and Bender (1973) is cited on a Russian wiki page, which I could not verify: \url{https://ru.ruwiki.ru/wiki/%D0%A3%D0%BB%D0%B8%D1%82%D0%B8%D0%B9%D1%81%D0%BA%D0%B8%D0%B9_%D1%8F%D0%B7%D1%8B%D0%BA}. Pronunciation notes: \url{https://www.pacificworlds.com/yap/land/lang4.cfm}.
\section{Religious and Community Resources}
None found beyond the Habele organization above.
\section{Digital and Audio Resources}
Pacific Worlds pages on the Ulithi language: \url{https://www.pacificworlds.com/yap/stories/lang5.cfm}, \url{https://www.pacificworlds.com/yap/sea/lang3.cfm}, \url{https://pacificworlds.com/yap/arrival/lang1.cfm}. I did not check them for audio.
\section{Curriculum Development Context}
The language is widely spoken, so a curriculum could build on the Habele dictionary and community knowledge. No existing school programme was found.
\end{document}